\documentclass[12pt,a4paper]{article}

% ---------------------------------------------------------
% PACKAGES
% ---------------------------------------------------------
\usepackage[utf8]{inputenc}
\usepackage[T1]{fontenc}
\usepackage{geometry}
\usepackage{enumitem}
\usepackage{longtable}
\usepackage{array}
\usepackage{booktabs}
\usepackage{hyperref}
\usepackage{xcolor}
\usepackage{listings}
\usepackage{float}
\usepackage{tikz}
\usetikzlibrary{arrows.meta, positioning, shapes.geometric}

\geometry{
    left=25mm,
    right=25mm,
    top=25mm,
    bottom=25mm
}

\hypersetup{
    colorlinks=true,
    linkcolor=blue,
    urlcolor=blue
}

% ---------------------------------------------------------
% CODE STYLE
% ---------------------------------------------------------
\definecolor{codegray}{rgb}{0.95,0.95,0.95}

\lstset{
    backgroundcolor=\color{codegray},
    basicstyle=\ttfamily\small,
    breaklines=true,
    frame=single,
    columns=fullflexible,
    showstringspaces=false
}

% ---------------------------------------------------------
% DOCUMENT
% ---------------------------------------------------------
\title{
    \textbf{Text Analyzer}\\
    \large MCA Final Project -- Software Architecture Specification
}

\author{}
\date{\today}

\begin{document}

\maketitle

\tableofcontents
\newpage

% =========================================================
% 1. PROJECT OVERVIEW
% =========================================================

\section{Project Overview}

\subsection{Project Name}

\textbf{Text Analyzer}

\subsection{Project Type}

Console-based modular application developed in the C programming language.

\subsection{Academic Level}

MCA Final Project.

\subsection{Primary Objective}

The objective of the project is to develop a robust and modular text
analysis application capable of accepting text from the user or from
text files and performing statistical, frequency, search, replacement,
and reporting operations.

The application must be designed using modular C programming practices
rather than placing all functionality inside a single source file.

\subsection{Primary Goals}

The system shall:

\begin{enumerate}
    \item Accept text directly from the user.
    \item Load text from external text files.
    \item Analyze text statistically.
    \item Analyze character and word frequencies.
    \item Search for words and phrases.
    \item Replace words or phrases.
    \item Save modified text.
    \item Generate analysis reports.
    \item Handle invalid input and file errors safely.
    \item Use dynamic memory where appropriate.
    \item Maintain clean separation between modules.
    \item Provide a testable and maintainable codebase.
\end{enumerate}

% =========================================================
% 2. TECHNOLOGY STACK
% =========================================================

\section{Technology Stack}

\begin{longtable}{p{0.30\textwidth} p{0.60\textwidth}}
\toprule
\textbf{Technology} & \textbf{Purpose} \\
\midrule
C & Core programming language \\
GCC & C compiler \\
Make & Build automation \\
Git & Version control \\
GitHub & Remote repository \\
Standard C Library & File, string, memory and character operations \\
\bottomrule
\end{longtable}

No external third-party C libraries shall be introduced unless there is
a clear technical requirement and the dependency is explicitly approved
before implementation.

% =========================================================
% 3. ARCHITECTURE
% =========================================================

\section{System Architecture}

The project shall follow a modular layered architecture.

\subsection{High-Level Architecture}

\begin{verbatim}
                         +----------------------+
                         |       main.c         |
                         | Application Entry    |
                         +----------+-----------+
                                    |
                                    v
                         +----------------------+
                         |        ui.c          |
                         | Menu / User Input    |
                         +----------+-----------+
                                    |
             +----------------------+----------------------+
             |                      |                      |
             v                      v                      v
      +-------------+        +-------------+        +-------------+
      | analyzer.c  |        |frequency.c  |        | search.c    |
      | Statistics  |        | Frequency   |        | Search/     |
      | Engine      |        | Analysis    |        | Replace     |
      +------+------+        +------+------+        +------+------+
             |                      |                      |
             +----------------------+----------------------+
                                    |
                                    v
                         +----------------------+
                         |  text_utils.c        |
                         | Common Text Helpers  |
                         +----------+-----------+
                                    |
                                    v
                         +----------------------+
                         | file_handler.c       |
                         | File I/O             |
                         +----------+-----------+
                                    |
                                    v
                         +----------------------+
                         |    Text Files        |
                         | data/ and reports/   |
                         +----------------------+
\end{verbatim}

% =========================================================
% 4. DIRECTORY STRUCTURE
% =========================================================

\section{Directory Structure}

The repository shall use the following structure:

\begin{verbatim}
text-analyzer/
|
+-- include/
|   +-- analyzer.h
|   +-- frequency.h
|   +-- file_handler.h
|   +-- search.h
|   +-- report.h
|   +-- text_utils.h
|   +-- ui.h
|   +-- common.h
|
+-- src/
|   +-- main.c
|   +-- analyzer.c
|   +-- frequency.c
|   +-- file_handler.c
|   +-- search.c
|   +-- report.c
|   +-- text_utils.c
|   +-- ui.c
|
+-- tests/
|   +-- test_analyzer.c
|   +-- test_frequency.c
|   +-- test_search.c
|   +-- test_text_utils.c
|
+-- data/
|   +-- sample.txt
|   +-- input.txt
|
+-- reports/
|   +-- .gitkeep
|
+-- docs/
|   +-- project_overview.md
|   +-- architecture.md
|   +-- test_cases.md
|   +-- user_manual.md
|
+-- build/
|   +-- .gitkeep
|
+-- Makefile
+-- README.md
+-- .gitignore
+-- LICENSE
\end{verbatim}

% =========================================================
% 5. MODULES
% =========================================================

\section{Module Architecture}

\subsection{main.c}

\textbf{Responsibility:}

The \texttt{main.c} file is the application entry point.

It shall:

\begin{enumerate}
    \item Initialize the application.
    \item Start the user interface.
    \item Control the application lifecycle.
    \item Perform final cleanup.
\end{enumerate}

The main file shall not contain business logic for text analysis.

\subsection{ui.c / ui.h}

Responsible for:

\begin{itemize}
    \item Main menu.
    \item Submenus.
    \item User input.
    \item Displaying results.
    \item Displaying errors.
    \item Displaying confirmations.
\end{itemize}

The UI module shall call appropriate service modules instead of
implementing analysis logic itself.

\subsection{analyzer.c / analyzer.h}

This module is the primary statistical analysis engine.

It shall calculate:

\begin{itemize}
    \item Total characters.
    \item Characters excluding whitespace.
    \item Total words.
    \item Total lines.
    \item Total sentences.
    \item Total paragraphs.
    \item Spaces.
    \item Tabs.
    \item Digits.
    \item Special characters.
    \item Uppercase letters.
    \item Lowercase letters.
    \item Vowels.
    \item Consonants.
    \item Longest word.
    \item Shortest word.
    \item Average word length.
    \item Average sentence length.
    \item Estimated reading time.
\end{itemize}

\subsection{frequency.c / frequency.h}

Responsible for:

\begin{itemize}
    \item Character frequency.
    \item Word frequency.
    \item Most frequent character.
    \item Most frequent word.
    \item Top-N words.
\end{itemize}

Frequency analysis shall be independent from the user interface.

\subsection{search.c / search.h}

Responsible for:

\begin{itemize}
    \item Word search.
    \item Phrase search.
    \item Occurrence counting.
    \item Case-sensitive search.
    \item Case-insensitive search.
    \item Word replacement.
    \item Phrase replacement.
\end{itemize}

\subsection{file\_handler.c / file\_handler.h}

Responsible for:

\begin{itemize}
    \item Opening files.
    \item Reading files.
    \item Writing files.
    \item Saving modified text.
    \item Checking file availability.
    \item Handling file-related errors.
\end{itemize}

\subsection{text\_utils.c / text\_utils.h}

Contains reusable text-processing helper functions.

Examples include:

\begin{itemize}
    \item Character classification.
    \item Vowel detection.
    \item Word-character detection.
    \item Sentence-ending detection.
    \item Whitespace handling.
    \item Case normalization.
    \item Word normalization.
    \item String trimming.
\end{itemize}

No duplicated helper logic should exist across analysis modules.

\subsection{report.c / report.h}

Responsible for generating human-readable analysis reports.

Reports shall contain:

\begin{itemize}
    \item Source information.
    \item General statistics.
    \item Character statistics.
    \item Word statistics.
    \item Frequency statistics.
    \item Search statistics when applicable.
\end{itemize}

Reports shall be stored inside the \texttt{reports/} directory.

% =========================================================
% 6. DATA STRUCTURES
% =========================================================

\section{Core Data Structures}

\subsection{TextStatistics}

The project shall define a structure similar to:

\begin{lstlisting}[language=C]
typedef struct {
    int characters;
    int characters_no_space;

    int words;
    int lines;
    int sentences;
    int paragraphs;

    int spaces;
    int tabs;
    int digits;
    int special_chars;

    int uppercase;
    int lowercase;

    int vowels;
    int consonants;

    double average_word_length;
    double average_sentence_length;
    double estimated_reading_time;
} TextStatistics;
\end{lstlisting}

\subsection{WordFrequency}

A word-frequency record may be represented as:

\begin{lstlisting}[language=C]
typedef struct {
    char *word;
    int count;
} WordFrequency;
\end{lstlisting}

\subsection{CharacterFrequency}

Character frequency may be represented as:

\begin{lstlisting}[language=C]
typedef struct {
    unsigned char character;
    int count;
} CharacterFrequency;
\end{lstlisting}

% =========================================================
% 7. COMMON DEFINITIONS
% =========================================================

\section{Common Definitions}

The \texttt{common.h} file shall contain shared constants and status
codes.

Example:

\begin{lstlisting}[language=C]
#ifndef COMMON_H
#define COMMON_H

#define MAX_FILENAME 256
#define MAX_WORD_LENGTH 100
#define DEFAULT_READING_SPEED 200

typedef enum {
    STATUS_SUCCESS = 0,
    STATUS_ERROR_FILE,
    STATUS_ERROR_MEMORY,
    STATUS_ERROR_INPUT,
    STATUS_ERROR_INVALID
} Status;

#endif
\end{lstlisting}

% =========================================================
% 8. PUBLIC INTERFACES
% =========================================================

\section{Module Interfaces}

\subsection{Analyzer Interface}

The analyzer module should expose functions similar to:

\begin{lstlisting}[language=C]
Status analyze_text(
    const char *text,
    TextStatistics *statistics
);

int count_words(const char *text);
int count_lines(const char *text);
int count_sentences(const char *text);
int count_paragraphs(const char *text);
\end{lstlisting}

\subsection{File Handler Interface}

\begin{lstlisting}[language=C]
Status load_file(
    const char *filename,
    char **content
);

Status save_file(
    const char *filename,
    const char *content
);
\end{lstlisting}

\subsection{Search Interface}

\begin{lstlisting}[language=C]
int search_word(
    const char *text,
    const char *word,
    int case_sensitive
);

int count_occurrences(
    const char *text,
    const char *pattern,
    int case_sensitive
);

Status replace_text(
    const char *text,
    const char *search,
    const char *replacement,
    char **result,
    int case_sensitive
);
\end{lstlisting}

The exact signatures may be adjusted during implementation if a
technically superior design is required, but unnecessary interface
changes should be avoided.

% =========================================================
% 9. FUNCTIONAL REQUIREMENTS
% =========================================================

\section{Functional Requirements}

\subsection{Text Input}

The system shall allow the user to enter text manually.

The system shall support multiline input.

\subsection{File Input}

The system shall allow the user to load a \texttt{.txt} file.

If the file cannot be opened, the application shall display a clear
error and return to the appropriate menu.

\subsection{Basic Analysis}

The system shall calculate all basic statistics defined in the
\texttt{TextStatistics} structure.

\subsection{Frequency Analysis}

The system shall calculate character and word frequencies.

Word frequency should be case-insensitive by default.

For example:

\begin{verbatim}
Computer
computer
COMPUTER
\end{verbatim}

may be normalized to:

\begin{verbatim}
computer
\end{verbatim}

before frequency calculation.

\subsection{Search}

The user shall be able to search for words or phrases.

Both case-sensitive and case-insensitive modes shall be supported.

\subsection{Replace}

The user shall be able to replace a word or phrase.

The application must not modify the original text until the user
explicitly confirms the operation.

\subsection{Report Generation}

The application shall generate a readable report in the
\texttt{reports/} directory.

% =========================================================
% 10. NON-FUNCTIONAL REQUIREMENTS
% =========================================================

\section{Non-Functional Requirements}

\begin{itemize}
    \item The code shall follow modular design principles.
    \item Functions should have single responsibilities.
    \item Global mutable state should be avoided.
    \item Dynamic memory must be released properly.
    \item File handles must always be closed.
    \item Invalid input must not crash the application.
    \item Compiler warnings should be minimized.
    \item Code should be portable across standard GCC environments.
    \item The application should remain usable for reasonably large
          text files.
    \item Business logic shall remain independent of UI code.
\end{itemize}

% =========================================================
% 11. MEMORY MANAGEMENT
% =========================================================

\section{Memory Management}

Text loaded from files may require dynamic memory allocation.

The lifecycle shall follow:

\begin{verbatim}
File
  |
  v
malloc / realloc
  |
  v
Text Buffer
  |
  +----> Analysis
  |
  +----> Search
  |
  +----> Frequency
  |
  +----> Report
  |
  v
free()
\end{verbatim}

Every successful dynamic allocation shall have a corresponding
deallocation.

Memory allocation failures shall be detected and reported.

% =========================================================
% 12. ERROR HANDLING
% =========================================================

\section{Error Handling}

The application shall handle at least the following conditions:

\begin{enumerate}
    \item File does not exist.
    \item File cannot be opened.
    \item File cannot be read.
    \item File cannot be written.
    \item Memory allocation failure.
    \item Empty input.
    \item Invalid menu selection.
    \item Empty search pattern.
    \item Empty replacement pattern where applicable.
    \item Invalid filename.
\end{enumerate}

Example error messages:

\begin{verbatim}
[ERROR] File not found.
[ERROR] Unable to open file.
[ERROR] Memory allocation failed.
[ERROR] Invalid menu choice.
[ERROR] Text cannot be empty.
\end{verbatim}

The application should recover from user errors without terminating
unexpectedly.

% =========================================================
% 13. USER INTERFACE
% =========================================================

\section{User Interface}

The application shall provide a menu-driven console interface.

\subsection{Main Menu}

\begin{verbatim}
========================================
             TEXT ANALYZER
========================================

1. Enter Text
2. Load Text File
3. View Current Text
4. Basic Statistics
5. Character Frequency
6. Word Frequency
7. Search Text
8. Replace Text
9. Generate Report
10. Save Text
0. Exit

Enter choice:
\end{verbatim}

\subsection{Analysis Menu}

\begin{verbatim}
========== ANALYSIS ==========

1. Character Count
2. Word Count
3. Line Count
4. Sentence Count
5. Paragraph Count
6. Letter Analysis
7. Digit Analysis
8. Special Character Analysis
9. Complete Statistics
0. Back
\end{verbatim}

\subsection{Frequency Menu}

\begin{verbatim}
======= FREQUENCY ANALYSIS =======

1. Character Frequency
2. Word Frequency
3. Most Frequent Character
4. Most Frequent Word
5. Top 10 Words
0. Back
\end{verbatim}

\subsection{Search Menu}

\begin{verbatim}
========== SEARCH ==========

1. Search Word
2. Search Phrase
3. Count Occurrences
4. Case-Sensitive Search
5. Case-Insensitive Search
6. Replace Word
0. Back
\end{verbatim}

% =========================================================
% 14. DATA FLOW
% =========================================================

\section{Data Flow}

\subsection{File Analysis Flow}

\begin{verbatim}
User
 |
 v
UI
 |
 v
File Handler
 |
 v
Text Buffer
 |
 +----------+----------+----------+
 |          |          |          |
 v          v          v          v
Analyzer Frequency   Search     Report
 |
 v
Statistics
\end{verbatim}

\subsection{Manual Input Flow}

\begin{verbatim}
User
 |
 v
UI
 |
 v
Input Buffer
 |
 v
Text Buffer
 |
 +----> Analyzer
 +----> Frequency
 +----> Search
 +----> Report
\end{verbatim}

% =========================================================
% 15. FREQUENCY ANALYSIS
% =========================================================

\section{Frequency Analysis}

\subsection{Character Frequency}

The application shall count occurrences of characters.

Example:

\begin{verbatim}
a : 25
b : 7
c : 14
d : 9
\end{verbatim}

\subsection{Word Frequency}

Words shall be normalized before comparison.

Example:

\begin{verbatim}
Computer computer COMPUTER
\end{verbatim}

may produce:

\begin{verbatim}
computer : 3
\end{verbatim}

\subsection{Top-N Words}

The application shall support displaying the most frequent words,
with Top 10 as the default user-facing option.

% =========================================================
% 16. SEARCH AND REPLACE
% =========================================================

\section{Search and Replace}

The system shall support:

\begin{itemize}
    \item Exact word search.
    \item Phrase search.
    \item Case-sensitive search.
    \item Case-insensitive search.
    \item Occurrence counting.
    \item Replacement.
\end{itemize}

Replacement should operate on a newly generated text buffer rather
than destructively modifying the original buffer.

Example:

\begin{verbatim}
Original:
I love C programming. C is powerful.

Search:
C

Replacement:
C language

Result:
I love C language programming. C language is powerful.
\end{verbatim}

The implementation must clearly define whether replacement operates on
whole words or arbitrary substrings. Whole-word replacement should be
preferred for word replacement.

% =========================================================
% 17. REPORT GENERATION
% =========================================================

\section{Report Generation}

A report shall contain at least:

\begin{verbatim}
========================================
        TEXT ANALYSIS REPORT
========================================

Source File:
Date/Time:

GENERAL STATISTICS
------------------
Characters:
Characters Without Spaces:
Words:
Lines:
Sentences:
Paragraphs:

CHARACTER ANALYSIS
------------------
Uppercase:
Lowercase:
Digits:
Spaces:
Tabs:
Special Characters:
Vowels:
Consonants:

WORD ANALYSIS
-------------
Longest Word:
Shortest Word:
Average Word Length:
Average Sentence Length:
Estimated Reading Time:

FREQUENCY ANALYSIS
------------------
Most Frequent Word:
Most Frequent Character:

========================================
\end{verbatim}

% =========================================================
% 18. TESTING
% =========================================================

\section{Testing Architecture}

Testing shall be separated from production code.

\subsection{Test Modules}

\begin{verbatim}
tests/
|
+-- test_analyzer.c
+-- test_frequency.c
+-- test_search.c
+-- test_text_utils.c
\end{verbatim}

\subsection{Test Categories}

The project shall test:

\begin{itemize}
    \item Normal input.
    \item Empty input.
    \item Single-word input.
    \item Multiline input.
    \item Punctuation.
    \item Numbers.
    \item Special characters.
    \item Mixed uppercase/lowercase text.
    \item Repeated words.
    \item Search failures.
    \item Replacement operations.
    \item File errors.
\end{itemize}

Example:

\begin{verbatim}
TEST: Word Count
Input: "Hello world"
Expected: 2
Actual: 2
Result: PASS
\end{verbatim}

% =========================================================
% 19. BUILD SYSTEM
% =========================================================

\section{Build System}

The project shall use a Makefile.

Required commands:

\begin{verbatim}
make
make run
make test
make clean
\end{verbatim}

The executable should be generated inside:

\begin{verbatim}
build/text_analyzer
\end{verbatim}

The build process should compile source files separately and then link
the resulting object files.

% =========================================================
% 20. COMPILER REQUIREMENTS
% =========================================================

\section{Compiler Requirements}

The project should be compiled using GCC with strong warning flags.

Recommended compilation flags:

\begin{verbatim}
-std=c11
-Wall
-Wextra
-Wpedantic
\end{verbatim}

Warnings should be treated seriously and fixed rather than ignored.

If the project later enables:

\begin{verbatim}
-Werror
\end{verbatim}

all existing warnings must first be resolved.

% =========================================================
% 21. GIT DEVELOPMENT STRATEGY
% =========================================================

\section{Git Development Strategy}

Development shall be performed incrementally.

\subsection{Phase 1}

\textbf{Project Setup}

\begin{itemize}
    \item Create directory structure.
    \item Create headers.
    \item Create source files.
    \item Create Makefile.
    \item Create README.
    \item Verify initial compilation.
\end{itemize}

Commit:

\begin{verbatim}
Phase 1: Initialize Text Analyzer project
\end{verbatim}

\subsection{Phase 2}

\textbf{Basic Text Input}

Implement:

\begin{itemize}
    \item Manual text input.
    \item Text buffer.
    \item View current text.
\end{itemize}

Commit:

\begin{verbatim}
Phase 2: Implement text input
\end{verbatim}

\subsection{Phase 3}

\textbf{Basic Statistics}

Implement:

\begin{itemize}
    \item Characters.
    \item Words.
    \item Lines.
    \item Sentences.
    \item Paragraphs.
\end{itemize}

Commit:

\begin{verbatim}
Phase 3: Implement basic text statistics
\end{verbatim}

\subsection{Phase 4}

\textbf{Advanced Statistics}

Implement:

\begin{itemize}
    \item Uppercase/lowercase.
    \item Digits.
    \item Special characters.
    \item Vowels/consonants.
    \item Longest/shortest word.
    \item Average lengths.
\end{itemize}

Commit:

\begin{verbatim}
Phase 4: Implement advanced text statistics
\end{verbatim}

\subsection{Phase 5}

\textbf{Character Frequency}

Commit:

\begin{verbatim}
Phase 5: Implement character frequency analysis
\end{verbatim}

\subsection{Phase 6}

\textbf{Word Frequency}

Commit:

\begin{verbatim}
Phase 6: Implement word frequency analysis
\end{verbatim}

\subsection{Phase 7}

\textbf{Search}

Commit:

\begin{verbatim}
Phase 7: Implement text search
\end{verbatim}

\subsection{Phase 8}

\textbf{Replace}

Commit:

\begin{verbatim}
Phase 8: Implement text replacement
\end{verbatim}

\subsection{Phase 9}

\textbf{File Handling}

Implement:

\begin{itemize}
    \item Load text files.
    \item Save text files.
    \item File validation.
    \item File error handling.
\end{itemize}

Commit:

\begin{verbatim}
Phase 9: Implement file handling
\end{verbatim}

\subsection{Phase 10}

\textbf{Report Generation}

Commit:

\begin{verbatim}
Phase 10: Implement analysis reports
\end{verbatim}

\subsection{Phase 11}

\textbf{Testing and Bug Fixing}

Commit:

\begin{verbatim}
Phase 11: Add tests and fix defects
\end{verbatim}

\subsection{Phase 12}

\textbf{Final Documentation}

Commit:

\begin{verbatim}
Phase 12: Complete documentation and final polish
\end{verbatim}

% =========================================================
% 22. DEVELOPMENT RULES
% =========================================================

\section{Development Rules}

The following rules are mandatory.

\begin{enumerate}

    \item Do not place the entire application inside
    \texttt{main.c}.

    \item Do not duplicate helper functions across modules.

    \item UI code shall not contain core analysis algorithms.

    \item File handling shall remain inside the file handler module.

    \item Text analysis shall remain independent of file I/O.

    \item Every dynamically allocated buffer shall have a clearly
    defined owner.

    \item Every allocated resource must eventually be released.

    \item Every opened file must be closed.

    \item Functions should perform one clearly defined responsibility.

    \item Avoid unnecessary global variables.

    \item Avoid unnecessary external dependencies.

    \item Do not modify unrelated files while implementing a feature.

    \item After every development phase, compile and test the project.

    \item Do not proceed to the next major phase if the current phase
    is broken.

    \item Do not silently ignore compiler warnings.

    \item Git commits should represent logical completed phases.

\end{enumerate}

% =========================================================
% 23. PERFORMANCE
% =========================================================

\section{Performance Considerations}

The analyzer should process a text file in a single or small number of
linear passes whenever practical.

For basic statistics, the target complexity should generally be:

\[
O(n)
\]

where \(n\) is the number of characters in the input.

Frequency and search algorithms may have higher complexity depending on
the chosen data structure, but unnecessary repeated full-text scans
should be avoided.

Memory usage should remain proportional to the size of the loaded text
and required analysis structures.

% =========================================================
% 24. SECURITY AND ROBUSTNESS
% =========================================================

\section{Robustness Requirements}

The implementation shall:

\begin{itemize}
    \item Avoid unsafe buffer operations where possible.
    \item Validate user input.
    \item Check all file operations.
    \item Check memory allocations.
    \item Prevent buffer overflows.
    \item Avoid dereferencing NULL pointers.
    \item Avoid use-after-free.
    \item Avoid double-free.
    \item Properly terminate strings.
\end{itemize}

Functions such as \texttt{gets()} shall never be used.

Safer alternatives such as \texttt{fgets()} shall be preferred.

% =========================================================
% 25. FUTURE EXTENSIONS
% =========================================================

\section{Future Extensions}

The architecture should allow future features without major redesign.

Possible future extensions include:

\begin{itemize}
    \item Readability scoring.
    \item Keyword extraction.
    \item Stop-word filtering.
    \item TF-IDF analysis.
    \item Basic sentiment analysis using a lexicon.
    \item Multiple file comparison.
    \item Text similarity.
    \item CSV report export.
    \item JSON report export.
    \item Graphical user interface.
\end{itemize}

These features are not part of the initial core implementation and
shall not be added prematurely.

% =========================================================
% 26. ACCEPTANCE CRITERIA
% =========================================================

\section{Acceptance Criteria}

The project shall be considered complete only when:

\begin{enumerate}

    \item The project builds successfully using the Makefile.

    \item The application starts without runtime errors.

    \item Text can be entered manually.

    \item Text files can be loaded.

    \item Basic statistics are calculated correctly.

    \item Advanced statistics are calculated correctly.

    \item Character frequency works correctly.

    \item Word frequency works correctly.

    \item Search works in both case-sensitive and
    case-insensitive modes.

    \item Replacement works without corrupting memory.

    \item Text can be saved.

    \item Reports can be generated.

    \item Invalid input is handled safely.

    \item File errors are handled safely.

    \item Dynamic memory is properly released.

    \item Test cases pass.

    \item The project compiles without significant warnings.

    \item README and documentation are complete.

    \item Git history contains logical development phases.

\end{enumerate}

% =========================================================
% 27. AGENT INSTRUCTIONS
% =========================================================

\section{Instructions for AI Coding Agent}

This document is the authoritative architecture specification for the
Text Analyzer project.

The coding agent shall follow these rules:

\begin{enumerate}

    \item Read and understand this architecture before modifying code.

    \item Implement the project incrementally according to the defined
    development phases.

    \item Do not implement all features in a single step.

    \item Before beginning a new phase, verify that the previous phase
    builds and works.

    \item Preserve the directory structure unless a strong technical
    reason requires modification.

    \item If a structural change becomes necessary, explain the reason
    before applying it.

    \item Do not introduce external dependencies without explicit
    approval.

    \item Prefer standard C and the standard C library.

    \item Keep modules independent and focused.

    \item Do not move business logic into the UI module.

    \item Do not duplicate code unnecessarily.

    \item Use meaningful function and variable names.

    \item Handle memory allocation and file operations safely.

    \item Compile with:
    
\begin{verbatim}
gcc -std=c11 -Wall -Wextra -Wpedantic
\end{verbatim}

    \item Run relevant tests after implementing each major feature.

    \item Do not claim a feature is complete until it has been
    compiled and tested.

    \item If an implementation decision is ambiguous, choose the
    simplest technically correct solution and document the decision.

    \item Do not add speculative features merely to make the project
    appear larger.

\end{enumerate}

% =========================================================
% 28. FINAL PROJECT STRUCTURE
% =========================================================

\section{Final Architecture Summary}

The final system shall follow:

\begin{verbatim}
                    TEXT ANALYZER
                          |
            +-------------+-------------+
            |                           |
        User Input                  File Input
            |                           |
            +-------------+-------------+
                          |
                          v
                    Text Buffer
                          |
          +---------------+---------------+
          |               |               |
          v               v               v
      Analyzer        Frequency        Search
          |               |               |
          +---------------+---------------+
                          |
                          v
                       Report
                          |
                          v
                    reports/*.txt
\end{verbatim}

The implementation shall prioritize correctness, modularity,
maintainability, memory safety, testability, and clear separation of
responsibilities.

\end{document}